% -------------------------------------------------------
% ASSIGNED TO: \textcolor{green}{Person 2}
% -------------------------------------------------------
% INSTRUCTION: Paste existing Section 3.5 introductory content here
% (the overview of EAs before the GA and NSGA-II subsections).
% Convert first-person to third person throughout.
% -------------------------------------------------------
 
Evolutionary algorithms (EAs) are a class of stochastic search and metaheuristic optimisation techniques 
inspired by the biological processes of natural evolution and selection. Unlike traditional algorithms that 
iterate on a single solution, EAs maintain a population of individuals, where each individual represents a 
potential solution to a problem (Yedjour et al., 2024). The population evolves over generations using three 
primary operators: selection, crossover, and mutation.  
\\\\
Rule extraction is a post-hoc explainability technique used to identify the underlying decision-making logic of 
opaque black-box models such as ANNs (Ghasemi et al., 2024; Yedjour et al., 2024). It is fundamentally a 
multi-objective optimisation problem requiring a simultaneous balance between fidelity to the black-box model, 
predictive accuracy on unseen data, and rule comprehensibility (Yedjour et al., 2024). Deterministic rule 
extraction methods such as decision trees impose fixed inductive biases on the search, for example, 
axis-aligned splits that limit depth and greedy coverage strategies, which constrain their ability to navigate 
the combinatorial complexity of high-dimensional genomic feature spaces or to optimise competing objectives 
simultaneously (Arrieta et al., 2020; Yedjour et al., 2024). EAs are well suited to this problem precisely 
because they maintain a diverse population of candidate rule sets, can be configured with flexible 
multi-objective fitness functions that reflect fidelity, accuracy, and comprehensibility simultaneously, and 
are capable of escaping local optima through mutation and crossover operators (Yedjour et al., 2024). Among EA 
variants, Genetic Algorithms and NSGA-II are selected for this project: GAs provide a single-objective baseline 
for rule extraction quality, while NSGA-II explicitly evolves a Pareto front of non-dominated rule sets, 
enabling a structured trade-off analysis between fidelity and comprehensibility that single-objective methods 
cannot produce (Yedjour et al., 2024). 
